\section*{Notation et domaines des constructions}
\phantomsection
\label{sec:ii-notacion}
\addcontentsline{toc}{section}{Notation et domaines des constructions}

Le tableau identifie les principaux symboles par leur domaine et par la
définition qui les introduit. Les indices, la typographie et le contexte
opératoire font partie de la notation. Les références renvoient aux
définitions effectives ; les occurrences anticipées conservent cette
même signification.

\begingroup
\small
\setlength{\tabcolsep}{4pt}
\renewcommand{\arraystretch}{1.12}
\begin{longtable}{@{}>{\raggedright\arraybackslash}p{0.19\linewidth}
                    >{\raggedright\arraybackslash}p{0.49\linewidth}
                    >{\raggedright\arraybackslash}p{0.27\linewidth}@{}}
\toprule
Symbole & Objet et signification & Définition et localisation\\
\midrule
\endfirsthead
\toprule
Symbole & Objet et signification & Définition et localisation\\
\midrule
\endhead
\bottomrule
\endfoot

\(\mathcal U_t\)
& Transition sur l'état à mémoire : sélection, transport et
  mise à jour des registres, dans cet ordre.
& \eqref{eq:actualizacion}.\\[3pt]

\(\mathcal R_{12}\), \(R_{\rm sgn}\)
& Registre de douze fenêtres avec chemin, orientation, retenue arithmétique et incidences ;
  lecteur signé à valeurs dans \(\ZZ^{12}\).
& \S\ref{subsec:k-registro-integral} ;
  composition \eqref{eq:k-lector-firmado}.\\[3pt]

\(A_m,C_m,V_m\)
& Trois dénombrements entiers par fenêtre : visites résiduelles, différence de
  traces orientées et incidence signée de frontière.
& \S\ref{subsec:registro-evaluacion-incidencial},
  définition des dénombrements.\\[3pt]

\(z_m\)
& Bloc décimal d'incidence dans \(\{0,\ldots,999\}\) :
  \(100[A_m]_{10}+10[C_m]_{10}+[V_m]_{10}\).
& \eqref{eq:registro-bloque-incidencial}.\\[3pt]

\(U\), \(U^{(r)}\)
& Vecteur entier signé à douze coordonnées et ses trois blocs
  harmoniques à quatre coordonnées.
& \eqref{eq:k-bloques-firmados} et
  \eqref{eq:k-u-firmado}.\\[3pt]

\(K\), \(K_j\)
& Registre entier \(K=\mathcal H_{12}U/4\), avec phase et orientation
  fixées ; prolongement de ses coordonnées par \(K_{j+12}=K_j\).
& \S\ref{subsec:k-hadamard}, définition de la transformation entière ;
  \S\ref{subsec:alpha-dominio}.\\[3pt]

\(\mathcal K_{\rm ph}\)
& Noyau de transfert sur \(\mathcal U_6\times C_9\),
  également appelé \(K_{\rm ph}\) lors de la description de la représentation à
  mémoire réduite.
& \S\ref{ii:tpk-arquitectura}, « Transfert réduit » ;
  retour \eqref{ii:tpk-renovacion}.\\[3pt]

\(A_j\), \(c_j\)
& Bloc de \(\alpha\) en base \(1000\) et quotient entier de la
  division euclidienne ; l'indice \(j\) parcourt le développement positionnel.
& \eqref{eq:alpha-recurrencia}.\\[3pt]

\(\mathbf z_m\), \(y_m\)
& Mémoire vectorielle accumulée dans \(\RR^{12}\) et coordonnée dans le
  repère mobile : \(y_m=S^{-m}\mathbf z_m\).
& \S\ref{sec:eta}, proposition de l'incrément ;
  mise à jour \eqref{eq:eta}.\\[3pt]

\(S\), \(R=S^{-1}\)
& Permutation cyclique de la base de mémoire,
  \(S\mathbf e_m=\mathbf e_{m+1}\) ; \(R\) désigne son inverse
  dans le développement de la résolvante.
& \S\ref{sec:eta} ;
  \S\ref{subsec:ii-memoria-resolvente}, paragraphe initial.\\[3pt]

\(\eta_m\)
& Incrément vectoriel produit par l'événement orienté :
  \(\eta_m=a_mS^{-1}\mathbf e_0\).
& \eqref{eq:events} et \eqref{eq:eta}.\\[3pt]

\(U:X\to X\), \(A\)
& Dans la réduction de mémoire, \(U\) est la mise à jour de l'état ;
  \(A:V_0\to V_0\) est le bloc conservé de son extension linéaire
  \(T\delta_x=\delta_{U(x)}\).
& \S\ref{subsec:memoria-resolvente}, « Élimination de mémoire » ;
  proposition~\ref{mem:prop-schur}.\\[3pt]

\(U:\mathcal H\to\mathcal H\)
& Réalisation isométrique du transport d'un tour,
  \(U^*U=I\), dans la compression observable.
& \S\ref{sec:extension}, « Compression observable » ;
  \eqref{eq:conservacion-ortogonal}.\\[3pt]

\(R=\ZZ/9\ZZ\)
& Anneau résiduel d'APP, d'idéal maximal
  \(\mathfrak m=(3)\).
& \S\ref{sec:nucleo}, « Réduction ternaire et radical multiplicatif ».\\[3pt]

\(R\), \(\eta\) locaux
& Dans le bloc central \(B_c=7I+2R\), \(R\) échange les deux
  coordonnées et \(R^2=I\). La rapidité de sa normalisation est
  \(\eta=\tfrac12\log(9/5)\).
& \S\ref{sec:extension}, « Développement local » :
  \(B_c/\sqrt{45}=\exp(\eta R)\).\\[3pt]

\(\Sigma_H\)
& Lecteur adimensionnel de norme et d’auto-échelle
  \(\Sigma_H=\varphi\alpha^{16}\) ; sa valuation décimale est
  \(\operatorname{ord}_{10}\Sigma_H=-34\).
& \eqref{eq:el-lector-determinantal} et \eqref{eq:el-decada}.\\[3pt]

\pagebreak[4]

\(\mathcal L_S\), \(\mathcal U_S\)
& Droite orientée d'action et base positive dans laquelle s'expriment
  ses sections.
& \S\ref{sec:action}, « Sections d'action et quotient adimensionnel
  de retour ».\\[3pt]

\(H_5\), \(\eta_{\rm ret}\)
& Coefficients adimensionnels d'action avant et après le retour ;
  \(\eta_{\rm ret}=H_5-(90/\pi)\mathcal C_\pi(\alpha)\).
& \eqref{eq:el-h5}--\eqref{eq:el-eta}.\\[3pt]

\(\hbar_{\rm pre}\), \(\hbar_{\rm ret}\), \(h_a\)
& Sections \(H_5\,10^{-34}\mathcal U_S\) et
  \(\eta_{\rm ret}\,10^{-34}\mathcal U_S\).
  Pour \(a\in\{\mathrm{pre},\mathrm{ret}\}\),
  \(h_a=2\pi\hbar_a\) exprime l'action par tour.
& \eqref{eq:el-secciones-accion} ;
  \S\ref{sec:action}, passage à l'action par tour.\\[3pt]

\(R_{\rm act}\), \(\delta_{\rm ret}\)
& Rapport multiplicatif \(H_5/\eta_{\rm ret}\) et différence
  additive \(H_5-\eta_{\rm ret}\) des coordonnées de retour.
& \eqref{eq:el-ratio}.\\[3pt]

\(A_{\deg}\), \(C^*_{\deg}\)
& Coordonnées angulaires \(1000\alpha\) et
  \(2(\eta_{\rm ret}+\alpha)\). Leurs coordonnées en radians sont
  \(x=\pi A_{\deg}/180\), \(y=\pi C^*_{\deg}/180\).
& \eqref{eq:ii-par-angular} ;
  \S\ref{sec:ii-angulos}, définition des coordinatisations.\\[3pt]

\(\gamma_{\HMT}\), \(\gamma\)
& Évaluation conjuguée de la chaîne orientée \(c_{12}^{+}\) en
  \(q_\pm\), avec composante angulaire \(\pi AC^*/180\), pour
  \(0<q_\pm<1\). L'incrément d'aire positif utilise
  \(|\gamma_{\HMT}|\).
& \S\ref{sec:ii-definicion-barbero} ;
  évaluation \eqref{eq:ii-definicion-barbero-evaluacion}.\\[3pt]

\(\eta_{\rm el}\), \(a_S,b_S\)
& Paramètre de forme \(\operatorname{artanh}(C^*/A)\) et
  demi-axes canoniques équilibrés de l'ellipse, avec
  \(a_Sb_S=2\hbar\).
& \S\ref{sec:ellipse}, définition initiale.\\[3pt]

\(R_{\rm el}\), \(\tau_{\rm el}\)
& Rayon modulaire adimensionnel \(\sqrt{b_S/a_S}\) et module
  \(\tau_{\rm el}=iR_{\rm el}^2\).
& \S\ref{sec:ellipse}, réciprocité de l'ellipse.\\[3pt]

\(K_a^i\)
& Courbure extrinsèque dans la convention normale de la réalisation
  géométrique ; elle intervient dans
  \(\mathcal A_a^i=\Gamma_a^i(E)+\gamma K_a^i\).
& \S\ref{sec:ii-costura-barbero}, définition de la connexion.\\[3pt]

\(L_*\), \(L\), \(\ell_P\)
& Référence longitudinale positive et longueur
  \(L=(5759/23040)\alpha_{\HMT}^{16}L_*\).
  La réalisation R1--R8 détermine \(\ell_P^2=L^2\) ;
  l’évaluation SI déclare \(L_*=1\,\mathrm m\).
& \eqref{g26:eq:regla-longitudinal-es} et
  \eqref{g26:eq:g-rigido-es}.\\[3pt]

\(R_X\), \(\Lambda_X\), \(E_L\)
& Lecture radiale positive, longueur conjuguée opératorielle et échelle
  énergétique \(E_L=\hbar_{\rm ret}c_{\rm int}/L\).
  Elles satisfont \(R_X\Lambda_X=L^2I\).
& \eqref{g26:eq:lecturas-conjugadas-es} et
  \eqref{g26:eq:productos-radiales-es}.\\[3pt]

\(R_{108}\), \(R\) d'horizon
& Rayon circulaire \(c_{\rm int}/\omega_{108}\) de la période
  dodécaphasique, égal à \(L\) pour l’horloge conjuguée ;
  \(R>0\) est la variable radiale de la réalisation
  d’horizon.
& \S\ref{sec:gravity}, définition initiale ;
  \eqref{eq:ii-horizonte-normalizado}.\\[3pt]

\(G_a\)
& Section gravitationnelle réciproque d’une action \(\hbar_a>0\) :
  \(G_a=c_{\rm int}^{3}L^{2}/\hbar_a\), avec \(R_{108}=L\),
  \(c_{\rm int}>0\) et coïncidence des rayons dans la convention
  \(r_g=Gm/c_{\rm int}^{2}\).
& \S\ref{sec:ii-definicion-gravedad};
  \eqref{eq:ii-definicion-gravedad-seccion} y
  \eqref{eq:ii-definicion-gravedad-area}.\\[3pt]

\(\mathbb G_{5,a}\)
& Opérateur \(G_aP^{\mathrm{ico}}_5\) sur
  \(V_{12}=\RR[A_5/C_5]\). Le projecteur orthogonal sélectionne
  le facteur direct \(V_5\) ; la section \(G_a\) est commune à ses cinq
  composantes.
& \eqref{ii:pent:eq:gravity-G5} ;
  théorème~\ref{ii:pent:thm:gravity-p5-coupling}.\\[3pt]

\(k_B\), \(\Theta_{{\rm clk},a}\)
& Transducteur linéaire positif
  \(k_B:\mathcal L_\Theta\to\mathcal L_E\) et section thermique
  positive, reliés par
  \(k_B\Theta_{{\rm clk},a}\log3=h_a/(108t_0)\).
  La section thermique et les bases déclarées déterminent sa
  coordonnée individuelle.
& \S\ref{sec:ii-definicion-boltzmann};
  \eqref{eq:ii-kb-aut-par-termico} y
  \eqref{eq:ii-kb-aut-coordenada}.\\

\end{longtable}
\endgroup

Dans le registre d’incidence, les fenêtres portent les indices \(1,\ldots,12\).
Dans la mise à jour de mémoire, elles sont numérotées de \(0\) à \(11\), et
leur prolongement utilise \(m\geq0\). Le décalage d’indice conserve
la même partition en fenêtres nonadiques. Dans les formules angulaires,
\(A,C^*\) conservent la coordinatisation déclarée dans chaque expression ; le quotient
\(C^*/A\) est invariant sous le changement commun entre degrés et radians.
